% =====================================================================
%  starter-faidah.tex — copy this to start a new FĀʾIDAH sheet.
%
%  A fāʾidah sheet is the notes from one lesson: no cover, no contents,
%  a header rule, and a title naming the lesson.
%
%  BUILD — from the kit's main folder:   bash tools/build.sh starter-faidah
%
%  THE ONE RULE THAT IS NOT TYPOGRAPHIC
%    Nothing here inserts content you did not type.
% =====================================================================
\documentclass[12pt]{article}
\newcommand{\ibfontpath}{fonts/}
\usepackage[faidah]{islamicbooklet}

% Optional — leave out for a plain sheet:
% \bookletdomain{www.example.org}   % header: "Brought to you by: www.example.org"
% \bookletlogo{my-logo.png}         % footer, bottom right

\begin{document}

% \faidahtitle{lesson code}{subject} — "Fawāʾid from:" is fixed.
\faidahtitle{L1}{The Subject of the Lesson}

At minute 00:00 of the class, the teacher asked the students about
something, and the answer was as follows.

\begin{enumerate}
  \item The first point.
  \item The second point.
\end{enumerate}

Taken from the āyah:

\begin{ayah}
  % \ar{...}  <- UNCOMMENT and type the Arabic here.
  %    Left commented so this starter builds as shipped: the Arabic
  %    faces carry no Latin glyphs, so a placeholder like "..." is a
  %    fatal "there are no glyphs" error, not a missing character.
  \en{...}
  \surah{Az-Zukhruf}{43:26-27}
\end{ayah}

\divider

% A plain Arabic paragraph — black by default, with the quotation itself
% coloured. One RTL run: isnād black, matn green, takhrīj black.
% \arpara{ ... \inhadith{«...»} ... }   <- UNCOMMENT and type the Arabic

The English underneath as its own paragraph, with the quotation itself
in \inhadith{green} and the grading black again.

\end{document}
